\newcommand{\EpisodeCount}{13}
\newcommand{\PrimaryScenarioCount}{325}
\newcommand{\DefaultDelay}{1}
\newcommand{\DefaultRefresh}{保守扣减}
\newcommand{\StudyConclusion}{\begin{itemize}\item 资金越大，按计划完成买卖越困难。参与率为 10\% 时，资金从 100 万元增至 1 亿元，当日完成率从 99.14\% 降至 86.19\%；1 亿元在前 30、60 分钟只完成 65.14\%、72.70\%。\item 大资金下，买入比卖出更难做完。1 亿元、10\% 参与率下，当日买入完成 80.55\%，卖出完成 97.47\%；13 个调仓日合计仍有 107 笔订单没有全部完成。\item 成交限制会改变当天收益。与全部按参考价成交相比，1 亿元、10\% 参与率下，13 个当日收益差的中位数为 -0.6769 个百分点。\end{itemize}}
\newcommand{\QualityNarrative}{原计划 14 个调仓日，每天 30 只股票，共应有 420 份单只股票的日行情。实际保留 13 天，共 1,678,487 条十档快照，其中连续竞价时段 1,618,662 条，能用于计算下单数量的 1,596,580 条。2026-07-01 的 30 只股票均无行情，已单独列出并排除。区间成交量和成交额共发现 130 条核对差异；连续竞价内，快照相隔超过 10 秒的情况有 21,173 次，最长 117 秒。}
\newcommand{\DailyQualityNarrative}{能与官方日线核对的 389 份日行情，日终累计成交量均在计算误差范围内一致。20 份的成交额有微小差别，最大为 0.40 元。000635.SZ 在 08-31 有 346 条快照，但没有官方日线，并查到当天停牌记录；因此用最初的参考价估值，另外标明价格没有更新，买入计划仍计入完成率。没有发现买卖档位顺序错误、连续竞价买价高于卖价，或数量达到字段上限的记录。}
\newcommand{\CompletionNarrative}{以 1 亿元为例，假设全部按参考价成交，三个时点自然都是 100\%。换成实际 L2 盘口，即使不限制参与率，30 分钟、60 分钟和当日也只能完成 93.38\%、95.91\%、98.43\%；参与率限制为 10\% 后，降为 65.14\%、72.70\%、86.19\%。这说明盘口数量本身已经构成限制，再限制每次下单数量，完成得更慢。分日期看，10\% 参与率下当日完成率为 82.24\% 至 90.51\%，13 天的中位数为 86.23\%。}
\newcommand{\StratumNarrative}{1 亿元、10\% 参与率下，成交较活跃的一组完成 100.00\%，成交较少的一组完成 71.00\%，成交一般的一组完成 90.31\%。股票按研究前的成交额分组，没有根据这次结果重新挑选。这只反映本次选定股票的差别，不能直接推广到全市场。}
\newcommand{\DurationNarrative}{1 亿元、10\% 参与率下，13 天共 390 笔买卖订单，做完 283 笔，未做完 107 笔。已经做完的订单，完成用时中位数为 14.35 个交易分钟；也就是其中一半不超过这个时间。未完成订单不参加这个中位数的计算。}
\newcommand{\ReturnNarrative}{1 亿元、10\% 参与率下，13 个调仓日的当日净收益中位数为 +0.4529\%。与全部按参考价成交相比，各日收益差从 -1.5595 到 +0.6308 个百分点不等。实际成交更少，有时反而会少亏一些，因此收益差为正不一定表示交易做得更好，还要看原计划是否完成。这些都是单日比较，不能相加或复利成三个月收益。}
\newcommand{\ClosingNarrative}{1 亿元、10\% 参与率下，未完成订单中有 2 笔买单、1 笔卖单，在 14:57 前最后 10 秒内的有效盘口中，已处于相应的官方涨跌停价，且原始盘口没有对手盘。这是事后查看收盘前行情得到的标记，只能说明那个时点的情况，不能证明全天都封板；本次也没有继续模拟这些剩余数量下一周能否成交。}
\newcommand{\SensitivityNarrative}{同样使用 1 亿元、10\% 参与率和原来的买卖计划，改变两项假设后，当日完成率为：每张快照重新计算可见数量、延迟 1 毫秒时，为 86.94\%；继续扣除当天已用过的同价位数量、延迟 3 秒时，为 86.19\%；继续扣除当天已用过的同价位数量、延迟 10 秒时，为 85.65\%。这些延迟都是实验参数，尚未测量真实下单耗时。1 毫秒和 3 秒的结果接近，可能因为行情快照之间本来就隔着数秒，两种延迟最终用到了同一张后续盘口。因此不能据此认定真实交易中慢 3 秒也没关系。}
\newcommand{\ValidationNarrative}{本轮程序通过 72 项测试，364 组运行结果均通过现金、持仓与负债核对。桌面重新运行 5 种成交规则，得到的 41 张结果表与原结果逐项一致。公开数据保留 10,920 行订单指标。重画图表时还会核对：计划股数等于已成交加剩余股数，未完成订单不填写完成时间，各图表使用同一份原计划金额。}
\newcommand{\CapitalRows}{%
100 & 99.14 & 99.14 & 99.14 & 3 & -0.2399\\
300 & 98.97 & 98.97 & 99.14 & 3 & -0.3822\\
1,000 & 98.05 & 98.44 & 98.96 & 5 & -0.5101\\
3,000 & 88.38 & 93.57 & 98.08 & 16 & -0.4832\\
10,000 & 65.14 & 72.70 & 86.19 & 107 & -0.6769\\
}
\newcommand{\QualityRows}{%
07-01 & 0/30 & 0 & -- & -- & --\\
07-06 & 30/30 & 126,165 & 1 & 1 & 2,064\\
07-13 & 30/30 & 126,706 & 1 & 1 & 1,894\\
07-20 & 30/30 & 131,652 & 0 & 0 & 1,351\\
07-27 & 30/30 & 124,977 & 0 & 0 & 2,170\\
08-03 & 30/30 & 131,712 & 15 & 15 & 1,245\\
08-10 & 30/30 & 132,122 & 13 & 13 & 1,343\\
08-17 & 30/30 & 129,585 & 4 & 4 & 1,622\\
08-24 & 30/30 & 129,708 & 10 & 10 & 1,602\\
08-31 & 30/30 & 126,999 & 17 & 17 & 1,671\\
09-07 & 30/30 & 131,509 & 41 & 41 & 1,350\\
09-14 & 30/30 & 127,376 & 7 & 7 & 1,827\\
09-21 & 30/30 & 131,170 & 13 & 13 & 1,317\\
09-28 & 30/30 & 128,806 & 8 & 8 & 1,717\\
}
